% Section B.2: the answers to the parts of lectures/L02/appendix/c_exercises.md that are not code.
% Every testbench message quoted was produced by a design built to Chapters 1 to 5 and broken in
% exactly the way the part describes.

\section{Chapter 2: UART Framing and the Transmitter}
\label{sol:c2}
The modules themselves are code: \code{baud_gen} and \code{uart_tx} in the a) parts of
Exercises~\ref{c2:ex:1} and~\ref{c2:ex:2}, checked by \code{baud_gen_tb} and \code{uart_tx_tb};
their instantiation into \code{uart_top} in Exercise~\ref{c2:ex:2} f), checked by a clean analysis
now and by \code{uart_top_tb} in \chapref{5}; and the parity and second stop bit of
Exercise~\ref{c2:ex:3} a), checked by your own testbench case.

\solution{c2:ex:1}{\code{baud_gen}}
\ans{b} \code{div = 50000000 / (16 * 115200) = 27.13}, which rounds to \textbf{27}, the value in the
protocol's table (\secref{proto:baud}). One bit is then $16 \times 27 = 432$ clocks, 8.64\,µs,
which is 115740.7~baud, 0.47\,\% fast.

The divide is by \code{16 * baud} because the receiver cannot see where a bit begins, so it samples
the line sixteen times per bit, and the transmitter and the receiver share the one time base.
\code{uart_rx} (\chapref{3}) spends the sixteen ticks: it detects the start edge on one of them,
counts to the middle of the start bit to re-check it, and then counts sixteen per bit to land each
later sample in the middle of its bit. \code{uart_tx} simply holds each bit for sixteen ticks. A
generator dividing by \code{baud} alone would give the transmitter everything it needs and leave the
receiver nothing to find the middle of a bit with.

\ans{c} The testbench passes with \code{counter = div - 1}. The counter reloads to zero on the very
edge it reaches \code{div - 1}, so with a constant \code{div} it never passes the threshold, and
``has reached'' and ``is at least'' are the same test: both fire exactly when the counter is 3.

They part company only when the counter is already past the threshold, and a \code{div} of zero is
the extreme case, since \code{div - 1} is then $-1$. With \code{>=}, every counter value is at least
$-1$, so \code{tick} is high on every clock and the counter stays at zero. With \code{=}, the
counter is never $-1$, so it never reloads: \code{tick} never rises, and the counter climbs until
its next increment leaves its declared range. Simulated with \code{div} forced to zero, the
\code{>=} version ticked on 99 of the first 100 clocks, and the \code{=} version never ticked and
stopped after 65536 clocks with

\begin{console}
ghdl:error: bound check failure at baud_gen.vhd:22
\end{console}

\noindent so \code{>=} is the one that keeps the counter inside its range. (In hardware the counter
would wrap silently instead, and still never tick.) In this design a zero cannot actually reach the
port, since \code{div} is \code{natural range 1 to 65535}: that is why \code{uart_top} substitutes
1 for a zero \code{BAUD_DIV}, and the robustness of \code{>=} is a second line of defence. It also
covers the case the range does allow, a \code{div} lowered while the counter is above the new
threshold: \code{>=} reloads on the next edge, \code{=} counts on to the top of its range first.

\ans{d} Both \code{uart_tx} and \code{uart_rx} count ticks, and both are synchronous: each looks at
\code{baud_tick} only inside \code{rising_edge(clock)}. A tick decoded combinationally from the
registered counter is still exactly one clock wide, because the counter holds \code{div - 1} for
exactly one clock. A glitch on it between edges, while the counter's bits change, is invisible to
both consumers as long as it has settled by the next edge, which is ordinary timing closure. So a
glitchy tick would not break them.

A \emph{wider} tick would. Every clock on which \code{baud_tick} is high counts as one sixteenth of
a bit, so a tick two clocks wide advances \code{uart_tx}'s \code{ticks} counter twice per period:
bits come out half as long as they should, and \code{uart_rx} samples at the wrong points. The
registered output is the design that makes ``one clock wide'' true by construction. A glitch would
matter only where \code{tick} is used as a clock or crosses into another clock domain, and it is
neither here.

\solution{c2:ex:2}{\code{uart_tx}}
\ans{b} The testbench stops at the first data bit:

\begin{console}
uart_tx_tb: data bit 0 mismatch!
\end{console}

\noindent \code{0x53} is \code{0101_0011}, so least significant first the data bits go out as
\code{1, 1, 0, 0, 1, 0, 1, 0}. Packed most significant first they go out as
\code{0, 1, 0, 1, 0, 0, 1, 1}: the first bit on the line is \code{data(7)}, a \code{0}, where the
testbench expects \code{data(0)}, a \code{1}. That is the byte shifted the other way, \code{0xCA}.

\ans{c} The testbench passes. \code{0xA5} is \code{1010_0101}. Least significant first, the eight data
levels are \code{1, 0, 1, 0, 0, 1, 0, 1}; most significant first they are
\code{1, 0, 1, 0, 0, 1, 0, 1} as well, because the byte reads the same in both directions, so the two
packings put the same ten levels on the line and a check of the levels cannot tell them apart.

For the check to be able to fail, the byte must \textbf{not be a bit-reversal palindrome}: its bit
$i$ must differ from its bit $7 - i$ for at least one $i$. A palindrome fixes the top four bits by
the bottom four, so there are $2^4 = 16$ of them, and those 16 lack the property: \code{0x00},
\code{0x18}, \code{0x24}, \code{0x3C}, \code{0x42}, \code{0x5A}, \code{0x66}, \code{0x7E},
\code{0x81}, \code{0x99}, \code{0xA5}, \code{0xBD}, \code{0xC3}, \code{0xDB}, \code{0xE7} and
\code{0xFF}.

\code{uart_rx_tb} sends \code{0x53} and \code{0x2D}, neither a palindrome, and \code{uart_top_tb}
sends \code{0x53}, so each constant is good. A loopback still cannot catch the two halves reversed
\emph{together}, however good its constant: a transmitter that sends \code{data(7)} first and a
receiver that stores the first bit it gets in bit 7 undo each other exactly, and the byte that
comes back is the byte that went in. With both halves reversed, \code{uart_top_tb} reports that it
passed. With either half reversed alone it fails with

\begin{console}
uart_top_tb: received byte mismatch, expected 0x53, got 0xCA!
\end{console}

\noindent and the receiver reversed alone fails its own testbench with the same message from
\code{uart_rx_tb}. Each unit testbench compares its module with the protocol, a fixed wire order,
rather than with the other module, so each catches its own half, and together they catch the pair
the loopback cannot.

\ans{d} \code{busy} is a level that is true for the whole frame, so it answers ``may another byte be
loaded now?'', which is the feeder's whole question; it is also half of \code{STATUS}'s TX-idle bit.
\code{done} is a one-clock event at the moment the stop bit ends, so it can count frames or raise a
transmit-complete interrupt exactly once per frame, which a level can only do with an edge detector
of its own.

A feeder that loaded on \code{done} instead of on \code{not busy} would never send anything:
\code{done} only pulses at the end of a frame, and no frame starts until something is loaded, so the
first byte waits for an event that only a sent byte can produce. Built that way, \code{uart_top_tb}
reports

\begin{console}
uart_top_tb: timed out waiting for RX valid!
\end{console}

\noindent because the byte written to \code{TX_DATA} never leaves the FIFO. Even with a first load
forced some other way, a one-clock pulse is the wrong thing to gate on: miss it, because the FIFO
was empty in that one clock and a byte arrived a cycle later, and that byte waits for a frame that
will never end.

\ans{e} \code{frame <= '1' & data & '0'} puts the leftmost item in the highest bit, so

\begin{narrowtable}{@{}lcccccccccc@{}}
  \thead{\code{frame} index} & 9 & 8 & 7 & 6 & 5 & 4 & 3 & 2 & 1 & 0 \\
  \midrule
  \thead{Bit} & stop & \code{d7} & \code{d6} & \code{d5} & \code{d4} & \code{d3} & \code{d2} &
    \code{d1} & \code{d0} & start \\
  \thead{Value} & \code{1} & & & & & & & & & \code{0} \\
\end{narrowtable}

\noindent Read from index 0 upward, that is the start bit, then \code{d0}, \code{d1} and so on to
\code{d7}, then the stop bit: exactly the wire order of an 8N1 frame. The caller may change
\code{data} on the very next cycle because \code{frame} is a register of the transmitter's own,
loaded once, in \code{STATE_IDLE}, on the edge that sees \code{start}. From then on every bit on the
line comes from \code{frame}, and \code{data} is not read again until the next \code{start} in the
next idle period. In \code{uart_top} that is what lets the FIFO advance on the same edge the
transmitter loads.

\solution{c2:ex:3}{Parity and a second stop bit}
\ans{b} The stop field, which is now one or two bits long. (With the parity bit of part a), the
parity field is present or absent too, so the frame is 10, 11 or 12 bits; the start bit and the
eight data bits never change.) At a fixed baud rate the frame time is the number of bits over the
baud rate, so a second stop bit makes every frame one bit period longer, 11 instead of 10: at
115200~baud that is 95.5\,µs instead of 86.8\,µs, and the throughput falls by a tenth,
from 11520 to 10473 bytes per second.

\ans{c} \code{uart_tx_tb} binds positionally and names eight ports, so it binds the eight existing
ones and leaves the three new inputs, declared after them, at their defaults of \code{'0'}. With
\code{parity_en} and \code{two_stop} low the frame is the old 8N1 frame, bit for bit, and the
testbench passes unchanged. (Declared anywhere else, or without a default, a new input shifts or
breaks the positional binding, and the testbench no longer analyzes against the module.)

To pin down even parity, drive \code{parity_en = '1'} and \code{parity_odd = '0'}, send \code{0x53},
and after the eight data bits sample one more bit centre: expect \code{'0'}, because \code{0x53}
has four 1s, already an even number. Then sample one bit later and expect the stop bit's
\code{'1'}. The choice of a byte whose parity bit is \code{'0'} is deliberate: a transmitter that
forgot to insert the parity bit would put its stop bit, a \code{'1'}, in that position, and one that
computed odd parity would send a \code{'1'} too, so both fail. Add a second byte of odd weight, such
as \code{0x52} with three 1s, and expect a \code{'1'}, so that a parity bit stuck at \code{'0'}
fails as well.
